% main.tex - showcase of the DC beamer theme
\documentclass[aspectratio=169,10pt]{beamer}
\usepackage{polyglossia}
\setdefaultlanguage{portuguese} % Replace with your desired language
\setotherlanguages{english, french, spanish} % Add any additional languages you need
\usepackage{fontspec}
\usetheme[showlogo]{DC}


% -- Metadata -----------------------------------
\title[DC Beamer Theme]{%
  Beamer Theme for DC%
  \texorpdfstring{{\footnotesize \quad \textcolor{DCred}{v1.1.0}}}{ v1.1.0}%
}
\author[D.Carvalho]{Prof. Diego Carvalho, D.Sc.}
\subtitle{Um template \LaTeX\ para apresentações no Cefet/RJ, baseado no UU Beamer Theme}
\institute{DIGES - DIREN/DEPES/CCEPRO - DIPPG/PPPRO - DIPPG/PPCIC}
\date{\today}

\begin{document}
% ===========================================
%  TITLE SLIDES  (single-author, then multi-author)
% ===========================================

% -- Single-author title --
\begin{frame}[plain]
  \titlepage
\end{frame}

% -- Multi-author / conference title (no subtitle - demonstrates that
%    leaving it out is fine) --
\subtitle{}
\author[Alice et al.]{%
  Alice\inst{1} \and
  Bob\inst{1} \and
  Eve\inst{2}%
}
\addinstitute{1}{Dept.\ of Information and Computing Sciences, Utrecht University}
\addinstitute{2}{Faculty of EEMCS, Delft University of Technology}
\multiauthor
\venue{NeurIPS 2026}
\begin{frame}[plain]
  \titlepage
\end{frame}

% -- Reset to single-author for the rest of the demo --
\subtitle{Um template \LaTeX\ para apresentações no Cefet/RJ}
\author[D.Carvalho]{Prof. Diego Carvalho, D.Sc.}
\institute{Cefet/RJ}
\clearinstitutes
\singleauthor

% ===========================================
%  TABLE OF CONTENTS
% ===========================================
\begin{frame}{Outline}
  \tableofcontents
\end{frame}

% ===========================================
\section{Iniciando os trabalhos}
% ===========================================

\begin{frame}{Usando o tema}
  Coloque no preâmbulo e defina o seu metadado:

  \begin{itemize}
    \item \texttt{\textbackslash usetheme\{DC\}} --- usagem básica
    \item \texttt{\textbackslash usetheme[showlogo]\{DC\}} --- coloca o logo do Cefet/RJ nos slides
    \item \texttt{\textbackslash usetheme[nl]\{DC\}} --- Dutch logo variant
  \end{itemize}

  \vskip 6pt
  Para diversos autores, utilize o padrão:
  \texttt{\textbackslash author} / \texttt{\textbackslash institute}
  with \texttt{\textbackslash and} and \texttt{\textbackslash inst\{n\}},
  then activate with:
  \begin{itemize}
    \item \texttt{\textbackslash multiauthor}
    \item \texttt{\textbackslash venue\{NeurIPS 2026, Utrecht\}} (opcional)
  \end{itemize}
\end{frame}

\begin{frame}{Opções do tema}
  \begin{columns}[T]
    \begin{column}{0.48\textwidth}
      \begin{block}{Flags} \footnotesize
        Passe opções para \texttt{\textbackslash usetheme[...]}\texttt{\{DC\}}:
        \texttt{nl}, \texttt{showlogo}, \texttt{nosectionpage}.
      \end{block}
      \vskip 6pt
      \begin{block}{Slides especiais} \footnotesize
        \texttt{\textbackslash standoutframe\{...\}} para ênfase.\\
        \texttt{\textbackslash thankframe\{...\}\{...\}} para fechamentos.\\
        \texttt{\textbackslash sectionframe\{...\}} divisor manual. \\
        \texttt{\textbackslash closingframe \{...\}} slide com o logo do Cefet para fechamentos.
      \end{block}
    \end{column}
    \begin{column}{0.48\textwidth} \footnotesize
      \begin{block}{Conference Title}
        \texttt{\textbackslash multiauthor} ativa o layout para múltiplos autores.
        Use \texttt{\textbackslash addinstitute\{n\}\{Name\}} para listar as afiliações
        (numbered automatically). Call \texttt{\textbackslash clearinstitutes} para resetar.
      \end{block}
    \end{column}
  \end{columns}

  \vskip 10pt
  A flag \texttt{showlogo}, ativada nesse exemplo, coloca o logo do Cefet no canto superior direito dos slides.
\end{frame}

% Manual section divider demo (same style as automatic \AtBeginSection pages)
\sectionframe{Esse é uma ativação de \textbackslash sectionframe}

% ===========================================
\sectionimage{img/palette.jpg}
\section{Colour Palette}
% ===========================================

\begin{frame}{DC Brand Colours}
  \begin{columns}[T]
    \begin{column}{0.48\textwidth}
      \begin{alertblock}{Cores principais}
        \textcolor{DCyellow}{\rule{1em}{1em}}~Yellow \texttt{\#FFCD00}\\[3pt]
        \textcolor{DCred}{\rule{1em}{1em}}~Red \texttt{\#C00A35}\\[3pt]
        {\rule{1em}{1em}}~Black \texttt{\#000000}\\[3pt]
        \textcolor{white}{\rule{1em}{1em}}~White \texttt{\#FFFFFF}
      \end{alertblock}
    \end{column}
    \begin{column}{0.48\textwidth}
    \begin{exampleblock}{Secondary Colours}
        \begin{tabular}{@{}l@{\hspace{0.25em}} l@{\hspace{2em}} l@{\hspace{0.25em}} l@{}}
        \textcolor{DCcream}{\rule{1em}{1em}}   & Cream      &
        \textcolor{DCorange}{\rule{1em}{1em}}  & Orange \\[3pt]

        \textcolor{DCburgundy}{\rule{1em}{1em}} & Burgundy   &
        \textcolor{DCbrown}{\rule{1em}{1em}}    & Brown \\[3pt]

        \textcolor{DCgreen}{\rule{1em}{1em}}   & Green      &
        \textcolor{DCblue}{\rule{1em}{1em}}    & Blue \\[3pt]

        \textcolor{DCdarkblue}{\rule{1em}{1em}} & Dark blue  &
        \textcolor{DCpurple}{\rule{1em}{1em}}   & Purple
        \end{tabular}
    \end{exampleblock}
    \end{column}
  \end{columns}

  \vskip 8pt
  Azul do Cefet é uma cor primária. \alert{Vermelho} é usado para chamar atenção.
  As cores secundárias são para os gráficos.
\end{frame}

% ===========================================
\section{Tipografia \& Blocos}
% ===========================================

\begin{frame}{Tipografia}
  \begin{enumerate}
    \item \textbf{Merriweather} (serif) para títulos e aberturas
    \item \textsf{Open Sans} (sans-serif) para texto do corpo
    \item Fallback: Palatino + se o pacote de fontes não funcionar
  \end{enumerate}

  \vskip 10pt
  \begin{exampleblock}{Tip}
    Instale \texttt{texlive-fonts-extra} para um visual autêntico.
  \end{exampleblock}
\end{frame}

\begin{frame}{Lists}
  Itemize utiliza quadrados marcadores em azul-cefet em até três níveis:

  \begin{itemize}
    \item First-level items use \textcolor{DCblue}{yellow} squares
    \begin{itemize}
      \item Second-level items use \textcolor{DCorange}{orange}
      \begin{itemize}
        \item Third-level items use \textcolor{DCbrown}{brown}
      \end{itemize}
    \end{itemize}
    \item Items can contain \textbf{bold}, \emph{italic}, or \alert{alerted} text
    \item Keep bullets concise
  \end{itemize}

  \vskip 8pt
  Enumerate uses numbered yellow squares:
  \begin{enumerate}
    \item First item
    \item Second item
    \item Third item
  \end{enumerate}
\end{frame}

\begin{frame}{Mathematics}
  The theme supports mathematical typesetting \cite{tai1994mathematical}:

$$
\int_{a}^{b} f(x)\,dx \approx
\frac{h}{2}\left[f(x_0) + 2\sum_{i=1}^{n-1} f(x_i) + f(x_n)\right]
$$

  \begin{alertblock}{Note}
    This is a note block, intended to grab the reader's attention.
  \end{alertblock}
\end{frame}

% ===========================================
\section{Diagrams \& Code}
% ===========================================

\begin{frame}{Workflow Diagram}
  \begin{figure}
  \centering
  \begin{tikzpicture}[
    node distance=0.6cm and 1.0cm,
    box/.style={
      rectangle, rounded corners=3pt,
      minimum height=2em, minimum width=4.5em,
      text centered, font=\sffamily\footnotesize,
      draw=none
    },
    arr/.style={->, >=stealth, semithick, DCdarkgray},
    scale=0.95, transform shape
  ]
    \node[box, fill=DCyellow, text=DCblack] (data) {Data};
    \node[box, fill=DCorange!80, text=white, right=of data] (preprocess) {Pre-process};
    \node[box, fill=DCgreen!90, text=white, right=of preprocess] (model) {Model};
    \node[box, fill=DCblue!90, text=white, right=of model] (evaluate) {Evaluate};
    \node[box, fill=DCburgundy!90, text=white, right=of evaluate] (deploy) {Deploy};

    \draw[arr] (data) -- (preprocess);
    \draw[arr] (preprocess) -- (model);
    \draw[arr] (model) -- (evaluate);
    \draw[arr] (evaluate) -- (deploy);

    \draw[arr, DCred, dashed]
      (evaluate.south) -- ++(0,-0.7) -| node[below, pos=0.25, font=\tiny\sffamily, DCred] {iterate} (model.south);
  \end{tikzpicture}
  \caption{Example process diagram with a numbered figure caption.}
  \end{figure}

  \vskip 6pt
  TikZ diagrams integrate naturally with the DC colour palette.
  Use \texttt{DCyellow}, \texttt{DCgreen}, \texttt{DCblue}, etc.\ directly.
\end{frame}

\begin{frame}[fragile]{Code Listing}
The \texttt{listings} package is pre-configured with DC colours.
  Specify the language and write code verbatim:

  \begin{lstlisting}[language=Python]
def fibonacci(n):
    """Return the n-th Fibonacci number."""
    if n <= 1:
        return n
    a, b = 0, 1
    for _ in range(2, n + 1):
        a, b = b, a + b  # shift window
    return b

print(fibonacci(10))  # 55
  \end{lstlisting}
\end{frame}

% ===========================================
\section{Tables \& Figures}
% ===========================================

\begin{frame}{Tables}
  \begin{table}
    \centering
    \caption{Theme files overview}
    \begin{tabular}{ll}
      \hline
      \textbf{File} & \textbf{Purpose} \\
      \hline
      \texttt{beamerthemeDC.sty}      & Main theme, options, fonts \\
      \texttt{beamercolorthemeDC.sty}  & Official DC brand colours \\
      \texttt{beamerinnerthemeDC.sty}  & Title page, blocks, special frames \\
      \texttt{beamerouterthemeDC.sty}  & Frame title, footline, logo \\
      \hline
    \end{tabular}
  \end{table}
\end{frame}

\begin{frame}{Figure in a Column Layout}
  \begin{columns}[T]
    \begin{column}{0.46\textwidth}
      Pair figures with text using \texttt{columns}:
      \begin{itemize}
        \item Left column carries the narrative
        \item Right column holds the figure
        \item Caption is numbered automatically
      \end{itemize}

      \vskip 8pt
      Use \texttt{[T]} on \texttt{columns} to top-align both sides,
      and \texttt{\textbackslash centering} inside the figure.
    \end{column}
    \begin{column}{0.50\textwidth}
      \begin{figure}
        \centering
        \begin{tikzpicture}[scale=0.82]
         % bars
          \foreach \x/\h/\col in {
            0.4/1.2/DCblue,
            1.1/2.4/DCblue,
            1.8/1.8/DCblue,
            2.5/2.9/DCyellow,
            3.2/2.1/DCyellow,
            3.9/3.0/DCyellow}{
            \fill[\col!80] (\x,0) rectangle (\x+0.55,\h);
          }
          % axes
          \draw[DCdarkgray, thin] (0,0) -- (4.5,0) node[right,font=\tiny\sffamily]{$x$};
          \draw[DCdarkgray, thin] (0,0) -- (0,3.2) node[above,font=\tiny\sffamily]{$y$};
          % legend
          \fill[DCblue!80]  (0.3,3.05) rectangle (0.65,3.25);
          \node[right,font=\tiny\sffamily] at (0.7,3.15) {Group A};
          \fill[DCyellow!80] (2.0,3.05) rectangle (2.35,3.25);
          \node[right,font=\tiny\sffamily] at (2.4,3.15) {Group B};
        \end{tikzpicture}
        \caption{Sample bar chart with DC colours.}
      \end{figure}
    \end{column}
  \end{columns}
\end{frame}

% -- Standout --------------------------
\standoutframe{  {\rmfamily\mdseries\LARGE\color{DCwhite} Sharing science,}\vspace*{-0.2cm}\\{\opensans\fontseries{l}\LARGE\itshape\color{DCwhite} shaping tomorrow}} % official spacing is a bit different though

% full image


\begin{frame}[plain]
\fullbleedimage{img/maracana2.jpg}
\end{frame}


% ===========================================
\section{Closing}
% ===========================================

\begin{frame}{Summary}
  \begin{enumerate}
    \item Blue-driven design, just as the DC likes it
    \item Merriweather + Open Sans, 10pt body text by default
    \item Special frames: \texttt{\textbackslash standoutframe}, \texttt{\textbackslash thankframe}, \texttt{\textbackslash closingframe}, \texttt{\textbackslash sectionframe},
    \item Options: \texttt{nl}, \texttt{showlogo}, \texttt{nosectionpage}
    \item Multi-author slides: \texttt{\textbackslash addinstitute\{n\}\{Name\}} + \texttt{\textbackslash multiauthor}
    \item Syntax-highlighted code via \texttt{lstlisting} with DC colours
  \end{enumerate}

  \vskip 12pt
  {\small Source and documentation: \url{https://github.com/diegomcarvalho/DC_Beamer_Template}}
\end{frame}

% -- Closing options demo --------------------------------------------
% (1) Compact thank-you with title + contact line
\thankframe{Feedback?}{Please contribute \\ \url{https://github.com/diegomcarvalho/DC_Beamer_Template}}

% (2) Brand closing slide — just the DC logo and tagline
\closingframe

% (3) Brand closing slide with partner-logo area
%     Replace the \rule boxes below with \includegraphics for real logos.
\closingframe[%
  \color{DCwhite!30}%
  \rule{2.6cm}{1.2cm}\hspace{0.8cm}%
  \rule{2.6cm}{1.2cm}\hspace{0.8cm}%
  \rule{2.6cm}{1.2cm}%
]

% -- references -----------------------------
\begin{frame}[allowframebreaks]{References}
  The Beamer class~\cite{tantau2004} provides the foundation.
  DC brand guidelines~\cite{uubrand2025} define the colour palette,
  typography, and logo usage rules.

  \vskip 6pt
  \printbibliography
\end{frame}

\end{document}
